% !TeX root = closed-systems-from-comparison-completeness.tex
% ============================================================
\chapter*{Orientation to Part I}
\phantomsection
\label{part:intro-foundations}
% ============================================================
\begin{chapterorientation}
\orientationitem{Settled}{The introduction fixes the closed-world admissibility stance and forbids external primitives as sources of distinction.}
\orientationitem{Task}{Part~I supplies the finite-to-global and comparison-completeness machinery needed before quotient semantics can be stated.}
\orientationitem{Handoff}{The output is the canonical product and diagonal-redundancy package used in Part~II.}
\end{chapterorientation}

Part~I works entirely before smooth geometry, field content, Hilbert
space, or observer-side dynamics.  Its role is to make explicit what
closed comparison data can and cannot determine at the finite and
inverse-limit levels.

\Cref{ch:foundation} isolates the Boolean compactness boundary:
finite-stage admissibility does not by itself imply global
admissibility unless the exclusion tower avoids a finite covering
failure.  \Cref{ch:bottom} applies that finite-to-global mechanism to
primitive comparison worlds.  The central result of the part is that the
rectangular product structure and diagonal redundancy used later are
derived structural consequences, not background scaffolding.
